% !TEX root = ../Main.tex
\chapter{微元思想}

\section{从匀速直线运动出发}

匀速直线运动 $S=vt$。在速度－时间（$v$–$t$）图上，速度是一条水平线，它与 $t$ 轴围成一个\textbf{矩形}：

\begin{center}
\begin{tikzpicture}[>=Stealth, scale=1]
    \draw[->] (0,0) -- (0,2.2) node[above] {$v$};
    \draw[->] (0,0) -- (3.4,0) node[right] {$t$};
    \draw[thick, blue] (0,1.4) -- (2.6,1.4);
    \draw[blue] (2.6,1.4) -- (2.6,0);
    \fill[pattern=north east lines, pattern color=blue!40] (0,0) rectangle (2.6,1.4);
\end{tikzpicture}
\end{center}

这个矩形的面积 $v\cdot t$ 恰好就是位移 $S$。也就是说，$v$–$t$ 图线与 $t$ 轴围成的矩形面积，代表了位移。

\section{变速直线运动呢？}

如果是变速直线运动，速度不再是水平线。先用一个大矩形去估：

\begin{center}
\begin{tikzpicture}[>=Stealth, scale=1]
    \draw[->] (0,0) -- (0,2.2) node[above] {$v$};
    \draw[->] (0,0) -- (3.6,0) node[right] {$t$};
    \draw[thick, blue, domain=0:3, samples=60, smooth] plot (\x, {0.6 + 1.1*(1 - exp(-1.1*\x))});
    \draw[dashed] (0,1.4) -- (3,1.4); \draw[dashed] (3,0) -- (3,1.4);
    \node[below] at (3,0) {$t_0$};
\end{tikzpicture}
\end{center}

这个矩形面积和真正的位移 $S$ 有"一定的贴近程度"，但\textbf{误差较大}。把 $t_0$ 一分为二，用两个矮一点的矩形去贴：

\begin{center}
\begin{tikzpicture}[>=Stealth, scale=1]
    \draw[->] (0,0) -- (0,2.2) node[above] {$v$};
    \draw[->] (0,0) -- (3.6,0) node[right] {$t$};
    \draw[thick, blue, domain=0:3, samples=60, smooth] plot (\x, {0.6 + 1.1*(1 - exp(-1.1*\x))});
    \foreach \a/\b in {0/1.5,1.5/3}{
        \pgfmathsetmacro\hh{0.6 + 1.1*(1 - exp(-1.1*\b))}
        \draw[dashed] (\a,\hh) -- (\b,\hh) -- (\b,0);
        \fill[pattern=north east lines, pattern color=blue!25] (\a,0) rectangle (\b,\hh);
    }
    \node[below] at (1.5,0) {\footnotesize $\tfrac{1}{2}t_0$};
    \node[below] at (3,0) {\footnotesize $t_0$};
\end{tikzpicture}
\end{center}

再分 $t_0$，分成更多份，每份用一个矩形去贴——\textbf{近似度更好了}：

\begin{center}
\begin{tikzpicture}[>=Stealth, scale=1]
    \draw[->] (0,0) -- (0,2.2) node[above] {$v$};
    \draw[->] (0,0) -- (3.6,0) node[right] {$t$};
    \draw[thick, blue, domain=0:3, samples=60, smooth] plot (\x, {0.6 + 1.1*(1 - exp(-1.1*\x))});
    \foreach \a/\b in {0/0.75,0.75/1.5,1.5/2.25,2.25/3}{
        \pgfmathsetmacro\hh{0.6 + 1.1*(1 - exp(-1.1*\b))}
        \draw[dashed] (\a,\hh) -- (\b,\hh) -- (\b,0);
        \fill[pattern=north east lines, pattern color=blue!25] (\a,0) rectangle (\b,\hh);
    }
    \node[below] at (0.75,0) {\footnotesize $\tfrac{1}{4}t_0$};
    \node[below] at (3,0) {\footnotesize $t_0$};
\end{tikzpicture}
\end{center}

让 $t_0$ 被\textbf{无限分割}，每一小段时间 $\Delta t$ 内速度几乎不变，那一小条的面积就是 $v\,\Delta t$；把所有小条加起来：

\[
\sum v\,\Delta t=S.
\]

\begin{center}
\begin{tikzpicture}[>=Stealth, scale=1]
    \draw[->] (0,0) -- (0,2.2) node[above] {$v$};
    \draw[->] (0,0) -- (3.6,0) node[right] {$t$};
    \draw[thick, blue, domain=0:3, samples=60, smooth] plot (\x, {0.6 + 1.1*(1 - exp(-1.1*\x))});
    \foreach \i in {0,...,11}{
        \pgfmathsetmacro\a{\i*0.25}
        \pgfmathsetmacro\b{(\i+1)*0.25}
        \pgfmathsetmacro\hh{0.6 + 1.1*(1 - exp(-1.1*\b))}
        \fill[pattern=north east lines, pattern color=blue!30] (\a,0) rectangle (\b,\hh);
        \draw (\a,0) rectangle (\b,\hh);
    }
    \node[below] at (3,0) {\footnotesize $t_0$};
\end{tikzpicture}
\end{center}

\state{微元思想的结论}{
    对于任意变速运动，$v(t)$ 图线与 $t$ 轴围成的图形面积，即为这段时间内的位移／路程。
}
